\section{Продолжение меры с полукольца}

\begin{definition}
    Пусть $\mu_0: \mathcal{P} \to [0, + \infty]$ — счётно-аддитивная мера на полукольце $\mathcal{P}$.
    \[\mu_0^*(E) := \inf \sum\limits_{j=1}^{\infty} \mu_0(P_j),\]
    где инфимум берётся по всем покрывающим множество $E$ последовательностям $\{P_j\} \subset \mathcal{P}$.
\end{definition} 

\begin{theorem}
    $\mu_0^*$ является внешней мерой и $\mu_0^*|_{\mathcal{P}} = \mu_0$.
\end{theorem} 
\begin{proof}
    Пусть $P \in \mathcal{P}$.
    Покажем, что $\mu_0^*(P) = \mu_0(P)$.
    Очевидно, что $P$ можно покрыть $\{P, \emptyset, \dots, \emptyset\}$.

    С одной стороны, $\mu_0^*(P) \leqslant \mu_0(P)$.
    С другой стороны, счётно-аддитивная мера на $\mathcal{P}$ является счётно-полуаддитивной, поэтому для любого покрытия $P \subset \bigcup\limits_{j=1}^{\infty} P_j$, где $\{P_j\} \subset \mathcal{P}$ верно $\mu_0(P) \leqslant \sum\limits_{j=1}^{\infty} \mu_0(P_j)$, поэтому, взяв инфимум по всем покрытиям, получим
    \[\mu_0(P) \leqslant \mu_0^*(P) \implies \mu_0(P) = \mu_0^*(P), \ \forall P \in \mathcal{P}.\]
    
    Покажем, что $\mu_0^*$ — счётно-полуаддитивная мера на $2^X$.
    Пусть $E \subset \bigcup\limits_{i=1}^{\infty} E_i$, где $E \in 2^X$, $\{E_i\} \subset 2^X$.
    Если существует $i \in \N$: $\mu_0^*(E_i) = + \infty$, то $\mu_0^*(E) \leqslant \sum\limits_{i=1}^{\infty} \mu_0^*(E_i)$ очевидно, поэтому далее считаем, что $\mu_0^*(E_i) \leqslant + \infty$ для любого $i \in \N$.

    Фиксируем $\epsilon > 0$ и для любого $i \in \N$ выберем $\{P_{i,j}\}_{j=1}^{\infty} \subset \mathcal{P}$ так, что 
    \[\mu_0^*(E_i) + \frac{\epsilon}{2^i} \geqslant \sum\limits_{j=1}^{\infty} \mu_0 (P_{i,j}).\]
    Итого $E \subset \bigcup\limits_{i=1}^{\infty} \bigcup\limits_{j=1}^{\infty} P_{i,j}$, следовательно, по определению $\mu_0^*$:
    \[\mu_0^* (E) \leqslant \sum\limits_{i=1}^{\infty} \sum\limits_{j=1}^{\infty} \mu_0(P_{i,j}) \leqslant \sum\limits_{i=1}^{\infty} \mu_0^*(E_i) + \frac{\epsilon}{2^i} \leqslant \epsilon + \sum\limits_{i=1}^{\infty} \mu_0^*(E_i).\]
    Поскольку $\epsilon > 0$ выбран произвольно, то $\mu_0^*(E) \leqslant \sum\limits_{i=1}^{\infty} \mu_0^*(E_i)$.
\end{proof} 

\begin{theorem}
    Пусть $\mu_0: \mathcal{P} \to [0, +\infty]$ — счётно-аддитивная мера на полукольце $\mathcal{P}$.
    Тогда $\mu_0^*$ является продолжением $\mu_0$ c $\mathcal{P}$ на $\sigma$-алгебру $\mathfrak{M}_{\mu_0^*}$, то есть:
    \begin{enumerate}
        \item $\mathcal{P} \subset \mathfrak{M}_{\mu_0^*}$;
        \item $\mu_0^*$ — $\sigma$-аддитивная мера на $\mathfrak{M}_{\mu_0^*}$;
        \item $\mu_0^*\vert_{\mathcal{P}} = \mu_0$.
    \end{enumerate}
\end{theorem} 
\begin{proof}
    2 и 3 утверждения теоремы следуют из прошлых теорем. Докажем первое.

    Требуется доказать, что 
    \[\forall E \subset X \ \forall P \in \mathcal{P} \ \mu_0^*(E) \geqslant \mu_0^*(P \cap E) + \mu_0^*(E \setminus P)\]
    (равенство в другую сторону очевидно).

    \begin{enumerate}
        \item Предположим, что $E \in \mathcal{P}$.
        Тогда $E \cap P \in \mathcal{P}$, где $P \in \mathcal{P}$, и $E \setminus P = \bigsqcup\limits_{j=1}^{N} Q_j$, где $\{Q_j\}_{j=1}^N \subset \mathcal{P}$.
        \begin{multline*}
            \underbrace{\mu_0(E)}_{\mu_0^*(E)} = (1) = \mu_0(E \cap P) + \sum\limits_{j=1}^{N} \mu_0(Q_j) = \mu_0^*(E \cap P) + \sum\limits_{j=1}^{N} \mu_0^*(Q_j) \geqslant (2) \\
            (2) \geqslant \mu_0^*(E \cap P) + \mu_0^*\left(\bigsqcup\limits_{j=1}^{N} Q_j\right) = \mu_0^*(E \cap P) + \mu_0^*(E \setminus P).
        \end{multline*}
        (1) — поскольку $\mu_0$ — мера на полукольце.
        (2) — ещё не знаем, что $Q_j$ измеримо, поэтому можем записать только такое неравенство.

        \item Пусть $E \subset X$ произвольно.
        Если $\mu_0^*(E) = + \infty$, то доказывать нечего.
        Если же $\mu_0^*(E) < + \infty$, то по определению инфимума зафиксируем $\epsilon > 0$ и найдём $\{P_j\}_{j=1}^{\infty} \subset \mathcal{P}$:
        \begin{multline*}
            \underbrace{\mu_0^*(E) + \epsilon}_{\bigcup\limits_{j=1}^{\infty} P_j \supset E} \geqslant \sum\limits_{j=1}^{\infty} \underbrace{\mu_0^*(P_j)}_{\geqslant \mu_0^*(P_j \cap P) + \mu_0^*(P_j \setminus P)} \geqslant \sum\limits_{j=1}^{\infty} \mu_0^*(P_j \cap P) + \mu_0^*(P_j \setminus P) \geqslant (1) \\
            (1) \geqslant \mu_0^*\left(P \cap \bigcup\limits_{j=1}^{\infty} P_j\right) + \mu_0^* \left(\bigcup\limits_{j=1}^{\infty} P_j \setminus P\right) \geqslant (2) \geqslant \mu_0^*(P \cap E) + \mu_0^*(E \setminus P).
        \end{multline*}
        (1) — в силу счётной полуаддитивности $\mu_0^*$.
        (2) — в силу монотонности $\mu_0^*$.

        В итоге для любого $\epsilon > 0$ верно
        \[\mu_0^*(E) + \epsilon \geqslant \mu_0^*(P \cap E) + \mu_0^*(E \setminus P) \implies \mu_0^*(E) \geqslant \mu_0^*(E \cap P) + \mu_0^*(E \setminus P).\]
    \end{enumerate}
\end{proof} 

\noindent \textbf{Вопрос.} А что, если ещё раз провернуть эту конструкцию не для полукольца $\mathcal{P}$, а для $\mathfrak{M}_{\mu_0^*}$?

\noindent \textbf{Ответ.} Ничего нового не получится.
С одной стороны, поскольку $\mathcal{P} \subset \mathfrak{M}_{\mu_0^*}$, то инфимум по большему запасу не может быть больше исходного инфимума.
Обозначим $\mu_0^*$ — мера на $\mathcal{P}$, $\mu^*$ — мера на $\mathfrak{M}_{\mu_0^*}$, тогда $\mu^* (E) \leqslant \mu_0^*(E)$.

С другой стороны,
\[\mu_0^*(E) \leqslant \sum\limits_{j=1}^{\infty} \mu_0^*(E_j) = \sum\limits_{j=1}^{\infty} \mu(E_j)\]
— продолженная на $\sigma$-алгебру мера.

\noindent Если $E \subset \bigcup\limits_{j=1}^{\infty} E_j$, $\{E_j\} \subset \mathfrak{M}_{\mu_0^*}$, то, взяв инфимум по всем покрытиям $\{E_j\}$:
\[\mu_0^*(E) \leqslant \mu^*(E) \implies \mu_0^*(E) = \mu^*(E) \ \forall E \subset X \implies \mathfrak{M}_{\mu_0^*} = \mathfrak{M}_{\mu^*}.\]

\begin{theorem}
    Пусть $\mu$ — конечно-аддитивная мера (со значениями $[0, + \infty]$) на $\sigma$-алгебре $\mathfrak{M}$.
    Тогда $\mu$ счётно-аддитивна тогда и только тогда, если $\mu$ непрерывна снизу по включению, то есть если
    \[A_1 \subset A_2 \subset \dots \subset A_n \subset \dots, \ \{A_i\} \in \mathfrak{M},\]
    то
    \[\exists \lim\limits_{i \to \infty} \mu(A_i) = \mu\left(\bigcup\limits_{i=1}^{\infty} A_i\right).\]
\end{theorem} 
\begin{proof}\tab

    \noindent\underline{$\implies$}
    Пусть $\mu$ счётно-аддитивна.
    \[B_1 = A_1, B_2 = A_2 \setminus A_1, \dots, B_n = A_n \setminus \bigcup\limits_{i=1}^{n-1} A_i.\]
    \begin{multline*}
        \bigcup\limits_{i=1}^{\infty} A_i = \bigsqcup\limits_{i=1}^{\infty} B_i \implies \mu\left(\bigcup\limits_{i=1}^{\infty} A_i\right) = \mu\left(\bigsqcup\limits_{i=1}^{\infty} B_i\right) =\\= \sum\limits_{i=1}^{\infty} \mu(B_i) = \underbrace{\sum\limits_{i=1}^{n} \mu(B_i)}_{\mu(A_n)} + \sum\limits_{i=n+1}^{\infty} \mu(B_i).
    \end{multline*}
    Если $\mu\left(\bigcup\limits_{i=1}^{\infty} A_i\right) < +\infty$, то остаток сходящегося ряда стремится к нулю, поэтому
    \[\mu\left(\bigcup\limits_{i=1}^{\infty} A_i\right) = \lim\limits_{n \to \infty} \mu(A_n).\]
    Для бесконечных значений утверждение следует из определения расходящегося ряда.

    \noindent\underline{$\impliedby$}
    Пусть $A = \bigsqcup\limits_{i=1}^{\infty} A_i$,
    \[\mu(A) = \lim\limits_{n \to \infty} \mu(B_n) = \lim\limits_{n \to \infty} \sum\limits_{i=1}^{n} \mu(A_i) = \sum\limits_{i=1}^{\infty} \mu(A_i).\]
    \[B_1 = A_1, \dots, B_n = \bigsqcup\limits_{i=1}^{n} A_i \implies B_1 \subset B_2 \subset \dots \subset B_n \subset \dots.\]
\end{proof} 

% конец 3 лекции

% начало 4 лекции
\paragraph{4 лекция.}

\begin{definition}
    Конечно-аддитивная мера $\mu$ на полукольце $\mathcal{P}$ подмножеств множества $X$ называется \textit{$\sigma$-конечной}, если 
    \[X = \bigcup\limits_{n=1}^{\infty} P_n, \ P_n \in \mathcal{P} \ \forall n \in \N\]
    и при этом 
    \[\mu(P_n) < + \infty \ \forall n \in \N.\]
\end{definition} 

\begin{definition}
    Тройка $(X, \mathfrak{M}, \mu)$ называется \textit{измеримым пространством}, если:
    \begin{itemize}
        \item $X$ — абстрактное множество;
        \item $\mathfrak{M}$ — $\sigma$-алгебра подмножеств множества $X$;
        \item $\mu: \mathfrak{M} \to [0, +\infty]$ — счётно-аддитивная функция на $\mathfrak{M}$ (называется мерой на $\mathfrak{M}$).
    \end{itemize}
\end{definition} 

\begin{definition}
    Пусть $(X, \mathfrak{M}, \mu)$ — измеримое пространство.
    Тогда $\mu$ называется \textit{полной мерой}, если $\forall E: \ \mu(E) = 0$ и $\forall E' \subset E$ выполнено $E' \in \mathfrak{M}$.
\end{definition} 

\begin{theorem}[О единственности продолжения меры]
    Пусть $\mu^*$ — продолжение $\sigma$-аддитивной меры $\mu$ с $\mathcal{P}$ на $\sigma$-алгебру $\mathfrak{M}_{\mu^*}$, $\nu$ — продолжение $\mu$ с $\mathcal{P}$ на $\sigma$-алгебру $\mathfrak{M}_{\nu}$.
    Тогда:
    \begin{enumerate}
        \item \[\forall A \in \mathfrak{M}_{\mu^*} \cap \mathfrak{M}_{\nu}: \ \nu(A) \leqslant \mu^*(A).\]
        Если $\mu(A) < +\infty$, то $\nu(A) = \mu^*(A)$.
        \item $\mu$ — $\sigma$-конечна, то 
        \[\forall A \in \mathfrak{M}_{\mu^*} \cap \mathfrak{M}_{\nu}: \ \nu(A) = \mu^*(A).\]
    \end{enumerate}
\end{theorem} 
\begin{proof}\tab
    \begin{enumerate}
        \item Пусть $A \subset \bigcup\limits_{j=1}^{\infty} P_j$.
        В силу счётной-полуаддитивности $\nu$:
        \begin{equation}
            \nu(A) \leqslant \sum\limits_{j=1}^{\infty} \nu(P_j) = (1) = \sum\limits_{j=1}^{\infty} \mu(P_j) \label{eq:4.1}
        \end{equation}
        (1) — поскольку $\nu$ — продолжение.
        Возьмём инфимум в \eqref{eq:4.1} по всем возможным покрытиям $A$ последовательностями $\{P_j\}$. Тогда:
        \[\nu(A) \leqslant \inf\limits_{A \subset \bigcup\limits_{j=1}^{\infty} P_j} \sum\limits_{j=1}^{\infty} \mu(P_j) = \mu^*(A).\]

        \item Докажем, что для любого $P \in \mathcal{P}$: $\mu(P) < +\infty$ и для любого $A \in \mathfrak{M}_{\mu^*} \cap \mathfrak{M}_{\nu}$
        \[\mu^*(\underbrace{P \cap A}_{\in \mathfrak{M}_{\mu^*} \cap \mathfrak{M}_{\nu}}) = \nu(P \cap A).\]
        Предположим противное, то есть
        \[\nu(P \cap A) < \mu^*(P \cap A).\]
        $P \in \mathfrak{M}_{\mu^*}$, то есть $P$ измеримо по Каратеодори.
        Тогда 
        \[\mu^*(P) = \nu(P) = \nu(P \cap A) + \nu(P \setminus A) < \mu^*(P \cap A) + \mu^*(P \setminus A) = \mu^*(P)\]
        — противоречие.

        \item Если $\mu$ $\sigma$-конечна, либо если $\mu^*(A) < +\infty$, то можно покрыть $A \subset \bigcup\limits_{j=1}^{\infty} P_j$: $\mu(P_j) < +\infty \ \forall j \in \N$ и $\{P_j\} \subset \mathcal{P}$.
        Но $\mu^*(A \cap P_j) = \nu(A \cap P_j)$, следовательно,
        \[\sum\limits_{j=1}^{\infty} \underbrace{\mu^*(A \cap P_j)}_{\leqslant \mu^*(A)} = \sum\limits_{j=1}^{\infty} \nu(A \cap P_j) = (*)\]
        По теореме о дизъюнктном представлении в полукольце можно считать, что $P_j$ попарно не пересекаются, поэтому
        \[(*) = \nu(A) \implies \mu^*(A) \leqslant \nu(A),\]
        а обратное неравенство мы уже имеем.
    \end{enumerate}
\end{proof} 

\begin{lemma}
    Если $\mu^*$ — внешняя мера, построенная по счётно-аддитивной мере $\mu$ на полукольце $\mathcal{P}$.
    Пусть есть $E \subset X$: $\mu^*(E) < +\infty$, тогда
    \[\exists C: \ C = \bigcap\limits_{n=1}^{\infty} \bigcup\limits_{j=1}^{\infty} P_{n,j}, \ P_{n,j} \in \mathcal{P},\]
    такое, что:
    \[\mu^*(C) = \mu^*(E).\]
\end{lemma} 
\begin{proof}
    Так как $\mu^*(E) < + \infty$, то $\forall n \in \N$ $\exists \{P_{n,j}\} \subset \mathcal{P}:$
    \[\frac{1}{2^n} + \mu^*(E) \geqslant \sum\limits_{j=1}^{\infty} \mu(P_{n,j}).\]
    Тогда
    \[C^N = \bigcap\limits_{n=1}^{N} \bigcup\limits_{j=1}^{\infty} P_{n,j} \supset E \ \forall N \in \N \implies 
    C = \bigcap\limits_{n=1}^{\infty} \bigcup\limits_{j=1}^{\infty} P_{n,j} \supset E,\]
    значит, $C \in \mathfrak{M}_{\mu^*}$.
    Тогда 
    \[\mu^*(C) \leqslant \mu^*(C^N) \leqslant \mu^*\left(\bigcup\limits_{j=1}^{\infty} P_{N,j}\right) \leqslant (1) \leqslant \sum\limits_{j=1}^{\infty} \mu^*(P_{N,j}) \leqslant \mu^*(E) + \frac{1}{2^N}.\]
    (1) — в силу счётной полуаддитивности $\mu^*$.

    А поскольку $N$ произвольно, то $\mu^*(C) \leqslant \mu^*(E)$. Но $C \supset E$, поэтому в силу монотонности $\mu^*$:
    \[\mu^*(E) \leqslant \mu^*(C) \implies \mu^*(C) = \mu^*(E).\]
\end{proof} 

\begin{remark}
    Заметим, что $C \in \mathfrak{M}(\mathcal{P})$, где $\mathfrak{M}(\mathcal{P})$ — наименьшая $\sigma$-алгебра, содержащая $\mathcal{P}$.
\end{remark}

\begin{theorem}
    Пусть $E \in \mathfrak{M}_{\mu^*}$ и $\mu^*(E) < + \infty$.
    Тогда $\exists B, C \in \mathfrak{M}(\mathcal{P}):$
    \[B \subset E \subset C\]
    и
    \[\mu^*(C \setminus B) = 0.\]
\end{theorem} 
\begin{proof}
    Берём $C$ из прошлой леммы.
    Определим \[e := C \setminus E\]
    — оно измеримо.
    Применим к $e$ предыдущую лемму, получим $\underbrace{\widetilde{e}}_{e \subset} \in \mathfrak{M}(\mathcal{P})$ и $\mu^*(\widetilde{e}) = 0$.
    Определим \[B := C \setminus \widetilde{e}.\]
    Тогда \[\mu^*(B) = \mu^*(E) = \mu^*(C).\]
\end{proof} 

\begin{theorem}
    Пусть $\mu$ — счётно-аддитивная и $\sigma$-конечная мера на $\mathcal{P}$, $\mathcal{P}$ — полукольцо подмножеств множества $X$.
    $\nu$ — полная мера на $\sigma$-алгебре $\mathfrak{M}_{\nu}$, $\nu$ — продолжение $\mu$.
    Тогда \[\mathfrak{M}_{\mu^*} \subset \mathfrak{M}_{\nu}.\]
\end{theorem} 
\begin{proof}
    По теореме о единственности представления меры:
    \[\nu\bigm|_{\mathfrak{M}(\mathcal{P})} = \mu^*\bigm|_{\mathfrak{M}(\mathcal{P})}.\]
    Фиксируем произвольное $e$: $\mu^*(e) = 0$ и покажем, что $e \in \mathfrak{M}_{\nu}$ и $\nu(e) = 0$.
    Поскольку $e \in \mathfrak{\mu^*}$ в силу полноты $\mu^*$, то
    \[\exists \widetilde{e} \in \mathfrak{M}(\mathcal{P}): \ \mu^*(\widetilde{e}) = 0 \text{ и } e \subset \widetilde{e} 
    \implies 
    \nu(\widetilde{e}) = 0 \text{ и } \widetilde{e} \in \mathfrak{M}_{\nu},\]
    но $\nu$ полна, поэтому $e \in \mathfrak{M}_{\nu}$.

    Теперь пусть $A \in \mathfrak{M}_{\mu^*}$ и $\mu^*(A) < +\infty$. Покажем, что $A \in \mathfrak{M}_{\nu}$.
    По предыдущей лемме: 
    \[\exists C \in \mathfrak{M}(\mathcal{P}) \subset \mathfrak{M}_{\nu}: \
    \mu^*(C) = \mu^*(A).\]
    Но $A$ измеримо, значит,
    \[C \setminus A \in \mathfrak{M}_{\mu^*} \text{ и } \mu^*(C \setminus A) = \mu^*(C) - \mu^*(A) = 0.\]
    Из $e = C \setminus A$ также следует, что $A = C \setminus e$.
    При этом $C$ и $e$ лежат в $\mathfrak{M}_{\nu}$, поэтому и $A \in \mathfrak{M}_{\nu}$.

    В общем случае, если $A$ только лишь измеримо (то есть не обязательно $\mu^*(A) < + \infty$), то в силу $\sigma$-конечности $\mu$
    \[X = \bigcup\limits_{n=1}^{\infty} X_n, \ X_n \in \mathcal{P}, \ \mu(X_n) < +\infty.\]
    Тогда
    \[A = \bigcup\limits_{n=1}^{\infty} (A \cap X_n) \text{ и }
    \forall n \in \N \ \mu^*(A \cap X_n) < +\infty,
    \]
    следовательно,
    \[A \cap X_n \in \mathfrak{M}_{\nu} \ \forall n \in \N
    \implies
    A \in \mathfrak{M}_{\nu}.\]
\end{proof} 